% ch21 — Ray Tracing in Curved Spacetime
\chapter{Ray Tracing in Curved Spacetime}
\label{ch:ray-tracing}

The centrepiece of the visualisation pipeline is the backward ray tracer that
integrates null geodesics from an observer's camera through the curved
spacetime to build an image. This chapter develops the algorithm from
initial-condition setup through GPU-accelerated integration.

\section{Backward Ray Tracing}
\label{sec:rt-backward}

We explain the backward ray-tracing paradigm: launching rays from each image
pixel back into the spacetime, rather than forward from luminous sources,
and justify its efficiency for image synthesis.

\section{Observer Tetrad and Initial Conditions}
\label{sec:rt-tetrad}
\coderef{Camera.Model}{hs}
\coderef{Camera.RayGenerator}{hs}

We construct the orthonormal tetrad of a stationary or orbiting observer,
map pixel coordinates to initial four-momenta on the observer's sky, and
normalise to the null condition $g_{\mu\nu} k^\mu k^\nu = 0$.

\section{Geodesic Integration}
\label{sec:rt-integration}

We formulate the geodesic equation as an ODE system for $(x^\mu, k_\mu)$
and integrate with the adaptive RK45 stepper developed in Chapter~8,
discussing step-size control and error tolerances.

\section{Termination Conditions}
\label{sec:rt-termination}

We define the stopping criteria: capture by the event horizon, escape to a
celestial sphere, maximum affine parameter, and intersection with a disk
or other emitting surface.

\section{GPU Acceleration with WGSL and wgpu}
\label{sec:rt-gpu}
\coderef{eh-gpu}{geodesic.wgsl}

We describe the port of the geodesic integrator to a WebGPU compute shader,
covering workgroup layout, buffer management, and the performance gains
from massively parallel ray integration.

\section{Double-Single Arithmetic}
\label{sec:rt-double-single}

We implement double-single (DS) arithmetic that achieves near-double precision
on GPU hardware limited to 32-bit floating point, and quantify the accuracy
improvement for long-integration geodesics.

\section{Geodesic Deviation and Ray Bundles\starmark{}}
\label{sec:rt-deviation}

We derive the geodesic deviation equation and show how it can be integrated
alongside the fiducial geodesic to propagate a ray bundle --- a small beam of
neighbouring geodesics represented by a Jacobian matrix. The beam
cross-section encodes the solid angle subtended by each pixel, enabling
gravitational lensing magnification estimates, adaptive supersampling in
regions of high demagnification, and correct flux weighting for extended
emitting sources.

\paragraph{Key References.}
Backward ray tracing follows \cite{James2015,Chan2013}; the observer-tetrad
construction draws on \cite{MTW1973}; double-single arithmetic follows
\cite{Thall2006}; GPU geodesic integration draws on \cite{Verbraeck2021};
ray-bundle transport follows \cite{Pretorius2009}.
